% Options for packages loaded elsewhere
\PassOptionsToPackage{unicode}{hyperref}
\PassOptionsToPackage{hyphens}{url}
\documentclass[
  a4paper,
]{ltjsarticle}
\usepackage{xcolor}
\usepackage[margin=25mm]{geometry}
\usepackage{amsmath,amssymb}
\setcounter{secnumdepth}{-\maxdimen} % remove section numbering
\usepackage{iftex}
\ifPDFTeX
  \usepackage[T1]{fontenc}
  \usepackage[utf8]{inputenc}
  \usepackage{textcomp} % provide euro and other symbols
\else % if luatex or xetex
  \usepackage{unicode-math} % this also loads fontspec
  \defaultfontfeatures{Scale=MatchLowercase}
  \defaultfontfeatures[\rmfamily]{Ligatures=TeX,Scale=1}
\fi
\usepackage{lmodern}
\ifPDFTeX\else
  % xetex/luatex font selection
\fi
% Use upquote if available, for straight quotes in verbatim environments
\IfFileExists{upquote.sty}{\usepackage{upquote}}{}
\IfFileExists{microtype.sty}{% use microtype if available
  \usepackage[]{microtype}
  \UseMicrotypeSet[protrusion]{basicmath} % disable protrusion for tt fonts
}{}
\makeatletter
\@ifundefined{KOMAClassName}{% if non-KOMA class
  \IfFileExists{parskip.sty}{%
    \usepackage{parskip}
  }{% else
    \setlength{\parindent}{0pt}
    \setlength{\parskip}{6pt plus 2pt minus 1pt}}
}{% if KOMA class
  \KOMAoptions{parskip=half}}
\makeatother
\usepackage{longtable,booktabs,array}
\newcounter{none} % for unnumbered tables
\usepackage{calc} % for calculating minipage widths
% Correct order of tables after \paragraph or \subparagraph
\usepackage{etoolbox}
\makeatletter
\patchcmd\longtable{\par}{\if@noskipsec\mbox{}\fi\par}{}{}
\makeatother
% Allow footnotes in longtable head/foot
\IfFileExists{footnotehyper.sty}{\usepackage{footnotehyper}}{\usepackage{footnote}}
\makesavenoteenv{longtable}
\usepackage{graphicx}
\makeatletter
\newsavebox\pandoc@box
\newcommand*\pandocbounded[1]{% scales image to fit in text height/width
  \sbox\pandoc@box{#1}%
  \Gscale@div\@tempa{\textheight}{\dimexpr\ht\pandoc@box+\dp\pandoc@box\relax}%
  \Gscale@div\@tempb{\linewidth}{\wd\pandoc@box}%
  \ifdim\@tempb\p@<\@tempa\p@\let\@tempa\@tempb\fi% select the smaller of both
  \ifdim\@tempa\p@<\p@\scalebox{\@tempa}{\usebox\pandoc@box}%
  \else\usebox{\pandoc@box}%
  \fi%
}
% Set default figure placement to htbp
\def\fps@figure{htbp}
\makeatother
\setlength{\emergencystretch}{3em} % prevent overfull lines
\providecommand{\tightlist}{%
  \setlength{\itemsep}{0pt}\setlength{\parskip}{0pt}}
% Glyph map for lualatex (newunicodechar). Maps symbols that may lack font glyphs.
\usepackage{newunicodechar}
\usepackage{amssymb}
\newunicodechar{∎}{\ensuremath{\blacksquare}}
\newunicodechar{✓}{\ensuremath{\checkmark}}
\newunicodechar{✗}{\ensuremath{\times}}
\newunicodechar{⟹}{\ensuremath{\Longrightarrow}}
\newunicodechar{⟺}{\ensuremath{\Longleftrightarrow}}
\newunicodechar{↪}{\ensuremath{\hookrightarrow}}
\newunicodechar{⌊}{\ensuremath{\lfloor}}
\newunicodechar{⌋}{\ensuremath{\rfloor}}
\newunicodechar{≃}{\ensuremath{\simeq}}
\newunicodechar{≈}{\ensuremath{\approx}}
\newunicodechar{≤}{\ensuremath{\leq}}
\newunicodechar{≥}{\ensuremath{\geq}}
\newunicodechar{≠}{\ensuremath{\neq}}
\newunicodechar{→}{\ensuremath{\rightarrow}}
\newunicodechar{←}{\ensuremath{\leftarrow}}
\newunicodechar{↔}{\ensuremath{\leftrightarrow}}
\newunicodechar{⊥}{\ensuremath{\perp}}
\newunicodechar{√}{\ensuremath{\surd}}
\newunicodechar{∑}{\ensuremath{\sum}}
\newunicodechar{∏}{\ensuremath{\prod}}
\newunicodechar{∞}{\ensuremath{\infty}}
\newunicodechar{∈}{\ensuremath{\in}}
\newunicodechar{∀}{\ensuremath{\forall}}
\newunicodechar{∣}{\ensuremath{\mid}}
\newunicodechar{γ}{\ensuremath{\gamma}}
\newunicodechar{λ}{\ensuremath{\lambda}}
\newunicodechar{μ}{\ensuremath{\mu}}
\newunicodechar{ρ}{\ensuremath{\rho}}
\newunicodechar{σ}{\ensuremath{\sigma}}
\newunicodechar{φ}{\ensuremath{\varphi}}
\newunicodechar{ψ}{\ensuremath{\psi}}
\newunicodechar{θ}{\ensuremath{\theta}}
\newunicodechar{Δ}{\ensuremath{\Delta}}
\newunicodechar{Π}{\ensuremath{\Pi}}
\newunicodechar{Λ}{\ensuremath{\Lambda}}
\newunicodechar{τ}{\ensuremath{\tau}}
\newunicodechar{ε}{\ensuremath{\varepsilon}}
\newunicodechar{ω}{\ensuremath{\omega}}
\newunicodechar{π}{\ensuremath{\pi}}
\newunicodechar{ν}{\ensuremath{\nu}}
\newunicodechar{±}{\ensuremath{\pm}}
\newunicodechar{×}{\ensuremath{\times}}
\newunicodechar{·}{\ensuremath{\cdot}}
\usepackage{bookmark}
\IfFileExists{xurl.sty}{\usepackage{xurl}}{} % add URL line breaks if available
\urlstyle{same}
\hypersetup{
  pdftitle={ロック後の幾何------元の面と終端面が張る4次元（2つの非零主角）、面の乗り換えは長窓SVDの集約アーティファクトでない},
  pdfauthor={Noriaki Kihara (WF System Co., Ltd.), ORCID 0009-0004-6753-4020},
  hidelinks,
  pdfcreator={LaTeX via pandoc}}

\title{ロック後の幾何------元の面と終端面が張る4次元（2つの非零主角）、面の乗り換えは長窓SVDの集約アーティファクトでない}
\author{Noriaki Kihara (WF System Co., Ltd.), ORCID 0009-0004-6753-4020}
\date{2026-09-12}

\begin{document}
\maketitle

\textbf{シリーズ:} 自己無撞着な関係波閉鎖系の基礎（論文7／全10本）\\
\textbf{著者:} 木原 範昭（WF System Co., Ltd.）　\textbf{日付:}
2026-09-12\\
\textbf{Version DOI:} 10.5281/zenodo.22729108\\
\textbf{Concept DOI:} 10.5281/zenodo.22729107\\
\textbf{ORCID:} 0009-0004-6753-4020\\
\textbf{Zenodo:} https://zenodo.org/records/22729108\\

\begin{center}\rule{0.5\linewidth}{0.5pt}\end{center}

\subsection{要旨}\label{ux8981ux65e8}

インフレーション（論文6）とSSBロック（論文5）が状態空間で何をするかを幾何的に確定する。写像は実直交回転
\(z(t)=Q_t z_0\) で、状態はどの瞬間も2次元面（rank
2）に留まる（\([\Re z,\Im z]\) の特異値が全 step
で二値、\(z^{\mathsf T}z=0\) 保存で直交等長フレーム）。床では床平面
\(P_0=\mathrm{span}\{\Re z_0,\Im z_0\}\) が \(K\) 不変で状態は \(P_0\)
に留まる。SSB で床を離れると、ロック面 \(P_{\rm term}\) は \(P_0\)
から傾き、両面の主角度は \(N\)
に依らず\textbf{二つとも非零}（\(N=6{:}27.1^\circ,30.8^\circ\)／\(N=40{:}18.1^\circ,20.4^\circ\)）で二面が張る次元は
\(4\)。\(P_0\) と \(P_{\rm term}\)
は2つの非零主角で4次元を張り、同じ床から終端の向きはシード依存で一意に定まらない。累積写像に特定の
\(SO(4)\) 回転構造は主張しない（実データで \(M_t\)
は4次元部分空間を保たないことを確認）。ロック面の床からの傾きは全
\(N=6\)--\(40\) で \(14.2^\circ\)--\(29.0^\circ\)（中央
\(\sim19.2^\circ\)）。遷移では第2方向が立ち上がり（transient
\(s_3/s_1\approx0.10\)--\(0.20\)＝実効
rank4）、ロック後は単一面へ再ロック（tail
\(s_3/s_1\approx10^{-5}\)--\(10^{-13}\)＝rank2）。短窓SVDテストにより、この「面の乗り換え」は長窓集約の見かけでなく\textbf{実在の面の傾き}であることを分離した。分類：導出済み幾何＋実測（アーティファクト排除）。

\begin{center}\rule{0.5\linewidth}{0.5pt}\end{center}

\subsection{1. 課題}\label{ux8ab2ux984c}

終点を「円環」と読むだけでは、この写像の本来の観察------\textbf{ある面から別の面へ乗り換える}------を捉えない。全
\(N\ge6\)
でこれを定量化し、見かけ（長窓SVDの集約）と実在の回転を分離する。

\subsection{2. 瞬時的構造------rank
2（単一面）}\label{ux77acux6642ux7684ux69cbux9020rank-2ux5358ux4e00ux9762}

写像は実直交回転 \(z(t)=Q_t z_0\)（\(Q_t\) 実直交、\(\Re z,\Im z\)
に同一作用）であり、状態はどの瞬間も2次元面に留まる（\([\Re z,\Im z]\in\mathbb R^{2\times M}\)
の特異値が全 step で二値のみ、\(z^{\mathsf T}z=0\)
保存で直交等長フレーム）。\textbf{Grassmann
運動としての定式化}：\(z=a+ib\) に対し \(z^{\mathsf T}z=0\)・\(\|z\|=1\)
から \(\|a\|^2=\|b\|^2=\tfrac12\)・\(a^{\mathsf T}b=0\) ゆえ
\(u=\sqrt2\,a,\ v=\sqrt2\,b\) は常に直交規格化された実2-frame
\((u,v)\in V_2(\mathbb R^M)\)。大域位相 \(z\mapsto e^{i\theta}z\)
は同じ面内で \((u,v)\)
を回すだけなので、観測対象は大域位相を除いた\textbf{2次元部分空間
\(P(t)=\mathrm{span}\{u,v\}\) そのもの}（Grassmann 多様体
\(\mathrm{Gr}(2,M)\) 上の運動）。射影
\(\Pi_t=u_tu_t^{\mathsf T}+v_tv_t^{\mathsf T}\) は各 step の実直交
\(Q_t=e^{\Delta\tau K_t}\) で厳密に

\[\Pi_{t+1}=Q_t\,\Pi_t\,Q_t^{\mathsf T}\]

と動く。床では \(Q_0P_0=P_0\) ゆえ
\(\Pi_{t+1}=\Pi_t\)（相対平衡・不動）、インフレで \(\Pi_t\neq\Pi_0\)
と平面が動き、飽和で \(P_{\rm term}\) へ再ロックする。§3--5
はこの平面運動 \(\Pi_0\to\Pi_t\to\Pi_{\rm term}\) を実測している。床では
\(P_0=\mathrm{span}\{\Re z_0,\Im z_0\}\) が \(K\)
不変（\(Ka=\mu b,\ Kb=-\mu a\)）で状態は \(P_0\)
に留まる（主角度ゼロ）。ロック後の末尾200 step の状態空間 SVD
第3/第1特異値比 \(s_3/s_1\) は
\(10^{-5}\)--\(10^{-13}\)＝単一面（rank2/\(S^1\)）。

\subsection{3. 初期面と終端面の張る次元------4D
span}\label{ux521dux671fux9762ux3068ux7d42ux7aefux9762ux306eux5f35ux308bux6b21ux51434d-span}

SSB で床を離れると、ロック面 \(P_{\rm term}\) は \(P_0\)
から傾く。元の面 \(P_0\) と終端面 \(P_{\rm term}\) の主角度は \(N\)
に依らず\textbf{二つとも非零}（\(N=6{:}27.1^\circ,30.8^\circ\)／\(N=40{:}18.1^\circ,20.4^\circ\)）で、二面が張る次元は
\(4\)。ロック面の床からの傾き（\(\arcsin\sqrt{H_\perp/H_{\rm end}}\)）は全
\(N=6\)--\(40\) で以下（\texttt{plane\_switch\_allN.csv} から抜粋）：

{\def\LTcaptype{none} % do not increment counter
\begin{longtable}[]{@{}llll@{}}
\toprule\noalign{}
\(N\) & 傾き{[}°{]} & transient \(s_3/s_1\) & tail \(s_3/s_1\) \\
\midrule\noalign{}
\endhead
\bottomrule\noalign{}
\endlastfoot
6 & 29.0 & 0.172 & \(3.8\times10^{-6}\) \\
8 & 20.3 & 0.135 & \(2.8\times10^{-4}\) \\
16 & 24.5 & 0.172 & \(1.1\times10^{-5}\) \\
22 & 15.7 & 0.130 & \(1.1\times10^{-13}\) \\
40 & 19.3 & 0.132 & \(4.5\times10^{-7}\) \\
64 & 21.7 & 0.021 & \(2.5\times10^{-4}\) \\
\end{longtable}
}

全 \(N\) で傾き \(14.2^\circ\)--\(29.0^\circ\)（中央
\(\sim19.2^\circ\)）、\(0^\circ\) でも \(90^\circ\)
でもない＝等方性が破れている。

\textbf{傾き・2主角・\(H_\perp/H\) は独立な測定量ではなく同一の
Grassmann 幾何。} 初期・終端フレーム \(U_0=(u_0,v_0),\ U_1=(u_1,v_1)\)
と \(P_0^\perp=I-U_0U_0^{\mathsf T}\)
について、主角の定義（\(U_0^{\mathsf T}U_1\) の特異値
\(=\cos\theta_i\)）から
\(\|P_0^\perp U_1\|_F^2=\sin^2\theta_1+\sin^2\theta_2\)。\(z=\tfrac1{\sqrt2}(u+iv)\)
ゆえ

\[\frac{H_\perp}{H}=\frac12\left(\sin^2\theta_1+\sin^2\theta_2\right)=\frac14\big\|\Pi_{\rm term}-\Pi_0\big\|_F^2,\qquad \alpha=\arcsin\sqrt{\frac{\sin^2\theta_1+\sin^2\theta_2}{2}}\]

（\(\Pi=UU^{\mathsf T}\) より
\(\|\Pi_1-\Pi_0\|_F^2=2(\sin^2\theta_1+\sin^2\theta_2)\)）。これは追加仮説でなく定義からの\textbf{厳密な幾何恒等式}で、\(N=6\)
の \((27.1^\circ,30.8^\circ)\to\alpha=29.0^\circ\)、\(N=40\) の
\((18.1^\circ,20.4^\circ)\to\alpha=19.3^\circ\)
が表と一致する。すなわち床からの傾き \(\alpha\) は初期面からの
\textbf{Grassmann 距離} \(\tfrac14\|\Pi_{\rm term}-\Pi_0\|_F^2\)
そのものである。

\begin{figure}
\centering
\pandocbounded{\includegraphics[keepaspectratio,alt={図1: 面の乗り換え（全N≥6）。左=ロック面の床からの傾き vs N、右=transient s3/s1（rank4/S³）→ tail s3/s1（rank2/S¹）。}]{figs/p7_postlock_geometry_ja_fig1.png}}
\caption{図1: 面の乗り換え（全N≥6）。左=ロック面の床からの傾き vs
N、右=transient s3/s1（rank4/S³）→ tail s3/s1（rank2/S¹）。}
\end{figure}

\textbf{図1} 面の乗り換えの \(N\) 依存。

\subsection{4.
元の面と終端面は4次元を張る}\label{ux5143ux306eux9762ux3068ux7d42ux7aefux9762ux306f4ux6b21ux5143ux3092ux5f35ux308b}

\(P_0\) と \(P_{\rm term}\)
は二つとも非零の主角度（\(N=6{:}27.1^\circ,30.8^\circ\)／\(N=40{:}18.1^\circ,20.4^\circ\)）で交わり
\(\dim(P_0+P_{\rm term})=4\)------終端面は元の面から2つの独立方向に傾く（この「4方向」は
span
の次元の意味に限る）。\textbf{同じ床からこの終端の向きはシード依存で一意に定まらない}（論文5）。累積写像
\(M_t=\prod_n\exp(\Delta\tau K(\varphi_n))\) に特定の \(SO(4)\)
回転構造は主張しない------実データで \(M_t\)
はこの4次元部分空間を保たない（\(M_tS=S\)
の不変性が破れる）ことを確認済みで、\(M_t\) は全 \(\mathbb R^M\)
の直交回転が \(P_0\) を4次元 span 内へ傾けるものである。

遷移期には第2回転面が立ち上がる（transient
\(s_3/s_1\approx0.10\)--\(0.20\)＝rank4）が、ロック後は単一面へ再ロック（tail
\(s_3/s_1\approx0\)＝rank2）。\(N=64\) は transient \(s_3/s_1=0.021\)
と小さく、大 \(N\) で第2面 transient が弱まる。この遷移 rank4
は一時的で、ロック後は rank2
に戻り、累積写像も4次元部分空間を保たない（§4）------\textbf{持続的な
\(S^3\) 構造は認められない}（検査済み、短窓SVD＋判別計算）。

論文6で確認した時変・非正規な接線コサイクルでは、最大増幅方向が時間とともに回転しながら増幅される。本論文で観測される
transient の追加方向は、その接線力学が瞬時 rank2 の平面 \(P(t)\)
を時間とともに傾ける幾何学的発現と読める。したがって transient の実効
rank4 は\textbf{瞬時状態の4次元化ではなく、移動する2平面 \(P(t)\)
が短時間に掃く局所 span} であり、飽和後には単一の終端面へ再び rank2
ロックする。すなわち論文6（接線空間で増幅方向が回転）と論文7（状態空間で2平面
\(P(t)\)
が傾いて別の2平面へ移る）は同一現象の接線側・状態側の表現である。\(H_\perp/H=\tfrac14\|\Pi_t-\Pi_0\|_F^2\)（§3
恒等式）を使えば、論文6の初期ランプ \(H_\perp/H\propto e^{2\gamma t}\)
は幾何学的に
\(\|\Pi_t-\Pi_0\|_F\propto e^{\gamma t}\)------\textbf{インフレーションは
Grassmann 空間 \(\mathrm{Gr}(2,M)\)
上での初期面からの指数的離脱}である。

\subsection{5.
アーティファクトの排除}\label{ux30a2ux30fcux30c6ux30a3ux30d5ux30a1ux30afux30c8ux306eux6392ux9664}

長窓の SVD
は、単一面が時間とともに回れば見かけ上高ランクに集約され得る。短窓SVDテスト（\texttt{artifact\_test\_shortwindow\_svd}）で、瞬時は
rank2 のまま元の面 \(P_0\) が2方向へ傾く（\(P_0\) と \(P_{\rm term}\)
が4次元を張る）ことを分離した------「面の乗り換え」は長窓集約の見かけでなく\textbf{実在の面の傾き}である。

\begin{figure}
\centering
\pandocbounded{\includegraphics[keepaspectratio,alt={図2: 短窓SVDテスト。瞬時rank2を保ちつつ元の面が2方向へ傾く（面の乗り換えは実在）。}]{figs/p7_postlock_geometry_ja_fig2.png}}
\caption{図2:
短窓SVDテスト。瞬時rank2を保ちつつ元の面が2方向へ傾く（面の乗り換えは実在）。}
\end{figure}

\textbf{図2} アーティファクト排除（短窓SVD）。

\subsection{6.
読み出しの但し書き}\label{ux8aadux307fux51faux3057ux306eux4f46ux3057ux66f8ux304d}

初期に複素平面が一枚のみのとき、その法線は縮退した直交補空間で観測不能（概念的）である。インフレーションが法線空間に特定の向きを立てて初めて、元の面と新しい面の相対位相（傾き）が読み出し可能になる。\textbf{読み出せるのは常に関係量（面どうしの相対位相）}であって、絶対的な法線方向は状態単独からは復元できない（\(Q_t\)
が向きを隠し、それを指す保存量もない）。

\subsection{7.
主張分類・未決}\label{ux4e3bux5f35ux5206ux985eux672aux6c7a}

\begin{itemize}
\tightlist
\item
  分類：導出済み幾何（実直交回転
  rank2）＋実測（主角度・傾き・4次元span・アーティファクト排除）＋判別計算（累積写像は4次元部分空間を保たず
  \(SO(4)\)
  ではないと確認）。判定：4次元span・2非零主角・相対配置の二重回転は保持、「力学が4次元上の
  \(SO(4)\) 二重回転」は棄却。
\item
  参考（別レジーム）：\(N=3,4\) の結晶終状態は面がマジック角
  \(\arccos\sqrt{2/3}=35.264^\circ\) で再凍結する特殊解（一般の
  \(N\ge6\) ロックは
  \(14\)--\(29^\circ\)）。結晶／ガラスの分岐は本シリーズの別課題。
\item
  未決：4次元 span の物理的同定（\(xyzt\) か \(xyztRQ\)
  か等）・ゲージ（目盛り）の型は本論文の範囲外（論文0 §未決）。
\end{itemize}

\subsection{関連研究（構造的一致・導出元でない）}\label{ux95a2ux9023ux7814ux7a76ux69cbux9020ux7684ux4e00ux81f4ux5c0eux51faux5143ux3067ux306aux3044}

\begin{itemize}
\tightlist
\item
  \(SO(4)\)
  の二重回転（不変平面2枚・独立回転角2つ）：\(P_0\to P_{\rm term}\)
  の\textbf{相対配置}の記述であって、累積写像そのものの性質ではない（§4
  の判別計算で \(M_tS\neq S\)）。
\item
  相対平衡（Marsden \& Ratiu）：床・ロックの剛体回転相対平衡。
\end{itemize}

\subsection{参考文献}\label{ux53c2ux8003ux6587ux732e}

\begin{enumerate}
\def\labelenumi{\arabic{enumi}.}
\tightlist
\item
  木原範昭「自己無撞着な関係波閉鎖系におけるインフレーション的急拡大の機構」Concept
  DOI 10.5281/zenodo.22112008.
\item
  J. E. Marsden and T. S. Ratiu, \emph{Introduction to Mechanics and
  Symmetry}, Springer (1999).
\end{enumerate}

\subsection{再現}\label{ux518dux73fe}

解析スクリプト・図・CSVは
\texttt{位相ロック全N確認\_相対平衡\_20260912/面の乗り換え全N\_20260912/}（\texttt{analyze\_plane\_switch\_allN\_20260912.py}、\texttt{artifact\_test\_shortwindow\_svd\_20260912.py}、\texttt{plane\_switch\_allN.csv}、\texttt{fig\_plane\_switch\_allN.png}、\texttt{fig\_artifact\_shortwindow\_svd.png}）。ロック済みデータは小
\(N\)・\(N=23\)--\(29\) が 500 step、未飽和 \(N\)（22,30--40）と
\(N=64\) が 2000 step。初期・終端面の主角度は
\texttt{N22\_2000step検証走行\_/}・\texttt{N40\_2000step検証走行\_/} の
step0/step2000 複素平面。早期の rank-3 監査は
\texttt{independent\_rank3\_geometry\_audit\_20260831/}。力学写像は正本
\texttt{run\_N3\_N40\_stage123\_v1.py}（SHA256
\texttt{1abf2353fee2e4f56f05e7a6f149fd086885136beb61ab571b48a56b09691567}）と同一・読出しのみ。

\end{document}
